\documentclass[11pt,a4paper]{scrartcl}

\usepackage{../../1.\space Vienanariai\space ir\space daugianariai/manostilius_lt}

\title{Vijeto teorema}
\author{Andrius Berniukevičius}

\begin{document}

\maketitle

\section{Kas yra Vijeto teorema?}

Išsprendus kvadratinę lygtį, jos sprendiniai (šaknys) ir pradinės lygties koeficientai yra glaudžiai susiję. Šį ryšį XVI amžiuje atrado prancūzų matematikas Fransua Vijetas (François Viète). 

\begin{definition}
\vocab{Redukuotoji kvadratinė lygtis} -- tai tokia kvadratinė lygtis, kurios pirmojo nario (prie $x^2$) koeficientas yra lygus vienetui ($a=1$). Tokia lygtis paprastai užrašoma pavidalu $x^2 + px + q = 0$.
\end{definition}

Jei turime pilnąją kvadratinę lygtį $ax^2 + bx + c = 0$, mes visada galime ją paversti redukuotąja, tiesiog padaliję abi lygties puses iš koeficiento $a$.

\begin{theorem}[Vijeto teorema pilnajai lygčiai]
Bet kokios kvadratinės lygties $ax^2 + bx + c = 0$ sprendinių sumai ir sandaugai galioja šios lygybės:
\[
\begin{aligned}
x_1 + x_2 &= -\frac{b}{a}, \\
x_1 \cdot x_2 &= \frac{c}{a}.
\end{aligned}
\]
\end{theorem}

\begin{proof}
Tarkime, kad lygtis turi du sprendinius. Pagal kvadratinės lygties šaknų formulę, tie sprendiniai yra:
\[ x_1 = \frac{-b - \sqrt{D}}{2a} \quad \text{ir} \quad x_2 = \frac{-b + \sqrt{D}}{2a}, \]
kur diskriminantas $D = b^2 - 4ac$.

Pirmiausia raskime šių sprendinių \textbf{sumą}:
\[
\begin{aligned}
x_1 + x_2 &= \frac{-b - \sqrt{D}}{2a} + \frac{-b + \sqrt{D}}{2a} \\
&= \frac{-b - \sqrt{D} - b + \sqrt{D}}{2a} \\
&= \frac{-2b}{2a} \\
&= -\frac{b}{a}.
\end{aligned}
\]
Sumos taisyklė įrodyta! Dabar raskime sprendinių \textbf{sandaugą}:
\[
\begin{aligned}
x_1 \cdot x_2 &= \frac{-b - \sqrt{D}}{2a} \cdot \frac{-b + \sqrt{D}}{2a} \\
&= \frac{(-b - \sqrt{D})(-b + \sqrt{D})}{4a^2}.
\end{aligned}
\]
Skaitiklyje pritaikome kvadratų skirtumo formulę $(x-y)(x+y) = x^2 - y^2$:
\[
\begin{aligned}
x_1 \cdot x_2 &= \frac{(-b)^2 - (\sqrt{D})^2}{4a^2} \\
&= \frac{b^2 - D}{4a^2}.
\end{aligned}
\]
Kadangi $D = b^2 - 4ac$, įsistatome šią išraišką vietoje $D$:
\[
\begin{aligned}
x_1 \cdot x_2 &= \frac{b^2 - (b^2 - 4ac)}{4a^2} \\
&= \frac{b^2 - b^2 + 4ac}{4a^2} \\
&= \frac{4ac}{4a^2} \\
&= \frac{c}{a}.
\end{aligned}
\]
Sandaugos taisyklė taip pat įrodyta[cite: 2].
\end{proof}

\begin{theorem}[Vijeto teorema redukuotajai lygčiai]
Jei redukuotoji kvadratinė lygtis $x^2 + px + q = 0$ turi sprendinius $x_1$ ir $x_2$, tai tų sprendinių suma yra lygi koeficientui prie $x$ su priešingu ženklu, o sprendinių sandauga yra lygi laisvajam nariui:
\[
\begin{aligned}
x_1 + x_2 &= -p, \\
x_1 \cdot x_2 &= q.
\end{aligned}
\]
\end{theorem}

\begin{proof}
Redukuotoji lygtis yra bendrosios kvadratinės lygties atvejis, kai
\[
a=1, \qquad b=p, \qquad c=q.
\]
Įstatę šias reikšmes į jau įrodytą bendrąją Vijeto teoremą, gauname:
\[
x_1+x_2=-\frac{b}{a}=-\frac{p}{1}=-p,
\qquad
x_1x_2=\frac{c}{a}=\frac{q}{1}=q.
\]
Teorema redukuotajai lygčiai įrodyta.
\end{proof}

\section{Vijeto teoremos taikymas}

Vijeto teorema dažniausiai naudojama atspėti lygties sprendinius neskaičiuojant diskriminanto, rasti lygties koeficientus, jei žinomi sprendiniai, arba patikrinti, ar lygtis išspręsta teisingai.

\subsection{Lygties sprendinių radimas ir spėjimas}

Dažnai naudinga tiesiog iš karto pasakyti lygties sprendinių sumą ir sandaugą, nesprendžiant pačios lygties.

\begin{example}
Neskaičiuodami lygties sprendinių, raskite jų sumą ir sandaugą:
\[ 2x^2 - 14x + 24 = 0 \]
\end{example}
\begin{solution}
Ši lygtis yra pilnoji, kur $a=2$, $b=-14$, $c=24$. Pagal bendrąją Vijeto teoremą:
\[
x_1 + x_2 = -\frac{b}{a} = -\frac{-14}{2} = 7
\]
\[
x_1 \cdot x_2 = \frac{c}{a} = \frac{24}{2} = 12
\]
\textbf{Atsakymas:} suma lygi $7$, sandauga lygi $12$.
\end{solution}

\begin{example}
Taikydami Vijeto teoremą, raskite lygties $x^2 - 5x + 6 = 0$ sprendinius.
\end{example}
\begin{solution}
Ši lygtis yra redukuotoji ($a=1$). Pagal Vijeto teoremą:
\[
\begin{aligned}
x_1 + x_2 &= -(-5) = 5, \\
x_1 \cdot x_2 &= 6.
\end{aligned}
\]
Mums reikia sugalvoti du skaičius, kuriuos sudauginę gautume $6$, o sudėję gautume $5$. 
Skaičių $6$ galime gauti dauginant $1 \cdot 6$ arba $2 \cdot 3$. 
Tikriname sumas:
$1 + 6 = 7$ (netinka),
$2 + 3 = 5$ (tinka).
Vadinasi, sprendiniai yra $2$ ir $3$.\\
\textbf{Atsakymas:} $x_1 = 2$, $x_2 = 3$.
\end{solution}

\begin{example}
Raskite lygties $x^2 + 2x - 15 = 0$ sprendinius.
\end{example}
\begin{solution}
Pagal Vijeto teoremą:
\[
\begin{aligned}
x_1 + x_2 &= -2, \\
x_1 \cdot x_2 &= -15.
\end{aligned}
\]
Kadangi sandauga neigiama ($-15$), sprendiniai turi būti skirtingų ženklų. 
Skaičių $15$ duoda sandaugos $1 \cdot 15$ ir $3 \cdot 5$. 
Kadangi suma neigiama ($-2$), didesnysis skaičius modulyje turi būti neigiamas.
Išbandome porą $3$ ir $-5$:
$3 \cdot (-5) = -15$ (tinka),
$3 + (-5) = -2$ (tinka).\\
\textbf{Atsakymas:} $x_1 = 3$, $x_2 = -5$.
\end{solution}

\subsection{Sprendinių sudarymas ir koeficientų radimas}

\begin{example}
Kvadratinės lygties $x^2 + px - 14 = 0$ vienas sprendinys yra $x_1 = 2$. Raskite antrąjį lygties sprendinį $x_2$ ir koeficientą $p$.
\end{example}
\begin{solution}
Pagal Vijeto teoremos sandaugos taisyklę žinome:
\[
\begin{aligned}
x_1 \cdot x_2 &= -14 \\
2 \cdot x_2 &= -14 \\
x_2 &= -7.
\end{aligned}
\]
Dabar pasinaudojame sumos taisykle, kad rastume koeficientą $p$:
\[
\begin{aligned}
x_1 + x_2 &= -p \\
2 + (-7) &= -p \\
-5 &= -p \\
p &= 5.
\end{aligned}
\]
\textbf{Atsakymas:} $x_2 = -7$, $p = 5$.
\end{solution}

\begin{example}
Sukurkite redukuotąją kvadratinę lygtį, kurios sprendiniai būtų $x_1 = 4$ ir $x_2 = -9$.
\end{example}
\begin{solution}
Naudojame Vijeto teoremą atvirkščiai. Raskime sumą ir sandaugą:
\[
\begin{aligned}
x_1 + x_2 &= 4 + (-9) = -5, \\
x_1 \cdot x_2 &= 4 \cdot (-9) = -36.
\end{aligned}
\]
Pagal teoremą žinome, kad $-p = -5$ (t. y. $p = 5$) ir $q = -36$.
Įstatome į formą $x^2 + px + q = 0$:
\[ x^2 + 5x - 36 = 0. \]
\textbf{Atsakymas:} $x^2 + 5x - 36 = 0$.
\end{solution}

\subsection{Simetrinių reiškinių skaičiavimas}

Vijeto teorema itin naudinga, kai norime apskaičiuoti reiškinius, sudarytus iš lygties sprendinių, bet pačių sprendinių žinoti (ar skaičiuoti) nenorime. Pavyzdžiui, sprendinių kvadratų sumą ($x_1^2 + x_2^2$) arba kubų sumą ($x_1^3 + x_2^3$). Mes šiuos reiškinius galime išreikšti per sprendinių sumą $(x_1+x_2)$ ir sandaugą $(x_1x_2)$, o jas jau žinome iš Vijeto teoremos.

\begin{remark}
Svarbios algebrinės tapatybės:
\[ x_1^2 + x_2^2 = (x_1 + x_2)^2 - 2x_1x_2 \]
\[ x_1^3 + x_2^3 = (x_1 + x_2)(x_1^2 - x_1x_2 + x_2^2) = (x_1 + x_2)\big((x_1 + x_2)^2 - 3x_1x_2\big) \]
\end{remark}

\begin{example}
Neskaičiuodami lygties $x^2 - 6x + 7 = 0$ sprendinių, raskite jų kvadratų sumą $x_1^2 + x_2^2$.
\end{example}
\begin{solution}
Pagal Vijeto teoremą žinome:
\[ x_1 + x_2 = 6 \]
\[ x_1 \cdot x_2 = 7 \]
Pasinaudojame tapatybe kvadratų sumai:
\[
\begin{aligned}
x_1^2 + x_2^2 &= (x_1 + x_2)^2 - 2x_1x_2 \\
&= 6^2 - 2 \cdot 7 \\
&= 36 - 14 = 22.
\end{aligned}
\]
\textbf{Atsakymas:} $22.$ 
\end{solution}

\begin{example}
Neskaičiuodami lygties $x^2 + 3x - 5 = 0$ sprendinių, raskite jų kubų sumą $x_1^3 + x_2^3$.
\end{example}
\begin{solution}
Pagal Vijeto teoremą:
\[ x_1 + x_2 = -3 \]
\[ x_1 \cdot x_2 = -5 \]
Pasinaudojame tapatybe kubų sumai:
\[
\begin{aligned}
x_1^3 + x_2^3 &= (x_1 + x_2)\big((x_1 + x_2)^2 - 3x_1x_2\big) \\
&= (-3)\big((-3)^2 - 3 \cdot (-5)\big) \\
&= (-3)(9 + 15) \\
&= (-3) \cdot 24 = -72.
\end{aligned}
\]
\textbf{Atsakymas:} $-72$ .
\end{solution}

\section{Uždaviniai}

\begin{problem}
Neskaičiuodami lygties sprendinių, nustatykite jų sumą ir sandaugą:
\begin{tasks}[label={(\arabic*)}, label-offset=0.9em](2)
    \task $x^2 - 9x + 20 = 0$
    \task $x^2 + 11x - 15 = 0$
    \task $3x^2 - 15x + 9 = 0$
    \task $-2x^2 - 8x + 10 = 0$
    \task $4x^2 + 12x - 7 = 0$
    \task $5x^2 - 20x + 15 = 0$
    \task $-3x^2 + 6x + 9 = 0$
    \task $7x^2 + 14x - 21 = 0$
    \task $2x^2 - 10x - 12 = 0$
    \task $-4x^2 - 4x + 8 = 0$
\end{tasks}
\end{problem}

\begin{problem}
Taikydami Vijeto teoremą, raskite šių kvadratinių lygčių sprendinius:
\begin{tasks}[label={(\arabic*)}, label-offset=0.9em](2)
    \task $x^2 - 7x + 10 = 0$
    \task $-x^2 + 8x - 12 = 0$
    \task $x^2 + 5x + 6 = 0$
    \task $-x^2 - 9x - 20 = 0$
    \task $x^2 - 3x - 10 = 0$
    \task $-x^2 - 4x + 21 = 0$
    \task $x^2 - x - 20 = 0$
    \task $-x^2 - x + 30 = 0$
\end{tasks}
\end{problem}

\begin{problem}
Žinomas vienas lygties sprendinys $x_1$. Raskite antrąjį sprendinį $x_2$ ir nežinomą koeficientą ($p$ arba $q$):
\begin{tasks}[label={(\arabic*)}, label-offset=0.9em](2)
    \task $x^2 + px + 12 = 0$, kai $x_1 = 3$
    \task $-x^2 - 6x + q = 0$, kai $x_1 = 4$
    \task $x^2 + px - 24 = 0$, kai $x_1 = -6$
    \task $-x^2 + 8x + q = 0$, kai $x_1 = -5$
    \task $-x^2 + 10x + q = 0$, kai $x_1 = 2$
    \task $x^2 - px - 18 = 0$, kai $x_1 = 6$
\end{tasks}
\end{problem}

\begin{problem}
Parašykite redukuotąją kvadratinę lygtį pavidalu $x^2 + px + q = 0$, jei žinomi jos sprendiniai:
\begin{tasks}[label={(\arabic*)}, label-offset=0.9em](2)
    \task $x_1 = 2$, $x_2 = 5$
    \task $x_1 = -3$, $x_2 = 7$
    \task $x_1 = -4$, $x_2 = -6$
    \task $x_1 = 1$, $x_2 = -8$
    \task $x_1 = 3$, $x_2 = -2$
    \task $x_1 = 0$, $x_2 = 4$
\end{tasks}
\end{problem}

\begin{problem}
Tarkime, kad $x_1$ ir $x_2$ yra lygties sprendiniai. Neskaičiuodami pačių sprendinių, raskite jų kvadratų sumą $x_1^2 + x_2^2$:
\begin{tasks}[label={(\arabic*)}, label-offset=0.9em](2)
    \task $-x^2 + 4x - 2 = 0$
    \task $x^2 + 5x - 3 = 0$
    \task $x^2 - 7x + 10 = 0$
    \task $-2x^2 + 8x - 4 = 0$
    \task $x^2 + 2x - 1 = 0$
    \task $-3x^2 + 6x - 1 = 0$
\end{tasks}
\end{problem}

\begin{problem}
Tarkime, kad $x_1<x_2$ yra lygties sprendiniai. Neskaičiuodami pačių sprendinių, raskite skirtumą $x_1^3-x_2^3$:
\begin{tasks}[label={(\arabic*)}, label-offset=0.9em](2)
    \task $x^2 - 5x + 6 = 0$
    \task $-x^2 - x + 6 = 0$
    \task $x^2 - 7x + 10 = 0$
    \task $-x^2 - 4x + 5 = 0$
    \task $x^2 - 6x + 8 = 0$
    \task $-x^2 + 2x + 8 = 0$
\end{tasks}
\end{problem}

\newpage
\answerssection

\begin{problem} \phantom{text}
\begin{tasks}[label={(\arabic*)}, label-offset=0.9em](2)
    \task suma $9$, sandauga $20$;
    \task suma $-11$, sandauga $-15$;
    \task suma $5$, sandauga $3$;
    \task suma $-4$, sandauga $-5$.
    \task suma $-3$, sandauga $-\frac{7}{4}$;
    \task suma $4$, sandauga $3$;
    \task suma $2$, sandauga $-3$;
    \task suma $-2$, sandauga $-3$;
    \task suma $5$, sandauga $-6$;
    \task suma $-1$, sandauga $-2$.
\end{tasks}
\end{problem}

\begin{problem} \phantom{text}
\begin{tasks}[label={(\arabic*)}, label-offset=0.9em](2)
    \task $x_1 = 2$, $x_2 = 5$;
    \task $x_1 = 2$, $x_2 = 6$;
    \task $x_1 = -2$, $x_2 = -3$;
    \task $x_1 = -4$, $x_2 = -5$.
    \task $x_1 = -2$, $x_2 = 5$;
    \task $x_1 = -7$, $x_2 = 3$;
    \task $x_1 = -4$, $x_2 = 5$;
    \task $x_1 = -6$, $x_2 = 5$.
\end{tasks}
\end{problem}

\begin{problem} \phantom{text}
\begin{tasks}[label={(\arabic*)}, label-offset=0.9em](2)
    \task $x_2 = 4$, $p = -7$;
    \task $x_2 = -10$, $q = 40$;
    \task $x_2 = 4$, $p = 2$;
    \task $x_2 = 13$, $q = 65$;
    \task $x_2 = 8$, $q = 32$;
    \task $x_2 = -3$, $p = 3$.
\end{tasks}
\end{problem}

\begin{problem} \phantom{text}
\begin{tasks}[label={(\arabic*)}, label-offset=0.9em](2)
    \task $x^2 - 7x + 10 = 0$;
    \task $x^2 - 4x - 21 = 0$;
    \task $x^2 + 10x + 24 = 0$;
    \task $x^2 + 7x - 8 = 0$;
    \task $x^2 - x - 6 = 0$;
    \task $x^2 - 4x = 0$.
\end{tasks}
\end{problem}

\begin{problem} \phantom{text}
\begin{tasks}[label={(\arabic*)}, label-offset=0.9em](2)
    \task $12$;
    \task $31$;
    \task $29$;
    \task $12$;
    \task $6$;
    \task $\frac{10}{3}$.
\end{tasks}
\end{problem}

\begin{problem} \phantom{text}
\begin{tasks}[label={(\arabic*)}, label-offset=0.9em](2)
    \task $-19$;
    \task $-35$.
    \task $-117$;
    \task $-126$;
    \task $-56$;
    \task $-72$.
\end{tasks}
\end{problem}

\end{document}